\clearpage
\SourceLeaf{338}
\pdfbookmark[2]{5. Gauss an Olbers, 7. Januar 1815}{parabolisch-005}
\begin{center}
[5.]\par
\medskip
\textsc{Gauss} an \textsc{Olbers}, Göttingen, 7. Januar 1815.
\end{center}
\noindent\hspace*{\parindent}.\ .\ .\ .\ .\ .\par
Ich habe Ihnen in meinem vorigen Brief die Näherungsformel für die Zwischenzeit
\[
t''-t=3mk\BareRoot{(r+r'')}\,.\BareRoot[10]{\left(1-\tfrac5{12}\frac{kk}{(r+r'')^2}\right)}\text{\begingroup\renewcommand{\thefootnote}{[*]}\footnotemark\endgroup}
\]
\HistoricalDisplayNoteText{Hier bezeichnet, abweichend von den vorstehenden Notizen, \(k\) die Sehne und \(\frac1{6m}\) die zuvor mit \(k\) bezeichnete Konstante, vgl. S. 330.]}
geschrieben und gesagt, dass sie in den meisten Fällen hinreichend genau sei. Ich hätte sagen sollen, dass sie in der Praxis stets mehr als genau genug wäre. Solange \(k<\tfrac12(r+r'')\), wirkt sich der Fehler noch nicht auf die 5\textsuperscript{te} Dezimalstelle aus. Es ist eine Freude zu sehen, wie bequem sich die Versuche für \(u\) damit durchführen lassen. Für jeden Versuch muss man nur 8-mal nachschlagen, 3-mal in den Logarithmentafeln, 4-mal in den Sinustafeln und einmal in meiner Hilfstafel für Logarithmen von Summen und Differenzen. Um möglichst wenig schreiben zu müssen, setze ich \(r\BareRoot{\frac{12}5}=s\) und \(r''\BareRoot{\frac{12}5}=s''\). Setzt man noch \(\frac{k}{s+s''}=\sin Q\), so ergibt sich
\[
\frac{t''-t}{m\BareRoot[4]{\frac{135}4}}=\frac{k\BareRoot{(s+s'')}}{\BareRoot[5]{\sec Q}}\mathpunct{\cdot}
\]

Die gesamte Vorrechnung besteht also nur noch im Folgenden.
\[
\begin{aligned}
\frac{bB}{u+c}&=\operatorname{tang}\theta,&\frac{b''B''}{u+c''}&=\operatorname{tang}\theta'',&\frac Au&=\operatorname{tang}\eta\\[5pt]
\frac{B\BareRoot{\frac{12}5}}{\sin\theta}&=s,&\frac{B''\BareRoot{\frac{12}5}}{\sin\theta''}&=s'',&\frac A{\sin\eta}&=\frac u{\cos\eta}=k.
\end{aligned}
\]

\SourceLeaf{339}\noindent\hspace*{\parindent}.\ .\ .\ .\ .\ .\par
Weil diese Versuche außerordentlich leicht sind, ist es ziemlich gleichgültig, mit welchem Wert von \(u\) man beginnt. Doch kann man auch hierbei einigermaßen methodisch vorgehen. Um nicht ganz im Dunkeln zu tappen, bestimme ich den ersten Wert von \(u\) auf folgende Art. Ich setze
\[
\frac{\BareRoot[6]{(A^4BB'')}}{\left(\frac{t''-t}{m\BareRoot{18}}\right)^{\frac23}}=\sin P,
\]
\[
u=-\tfrac16(c+c'')+\cos P\cdot\left(\frac{t''-t}{m\BareRoot{18}}\right)^{\frac23}\cdot\BareRoot[6]{(bb'')},
\]
was freilich nur eine höchst grobe Näherung sein kann. Der Grund für diese Formel ist folgender. Zunächst lasse ich \(\BareRoot[5]{\sec Q}\) außer Betracht. Ich setze
\[
\tfrac12(r+r'')=\BareRoot{\left\{\left(\frac{u+\frac12(c+c'')}{\BareRoot{(bb'')}}\right)^2+\BareRoot{(BB'')}^2\right\}}
\]
und, was freilich noch kühner ist,
\[
k^{\frac23}\left(\frac{r+r''}2\right)^{\frac13}=\BareRoot{\left\{\left(\frac{u+\frac16(c+c'')}{\BareRoot[3]{(\BareRoot{(bb'')})}}\right)^2+\left(A^{\frac23}(\BareRoot{(BB'')})^{\frac13}\right)^2\right\}}.
\]
Nachdem ich die Rechnung mit diesem ersten Wert von \(u\) durchgeführt habe, erhalte ich den Fehler bei \(\log\frac{t''-t}{m\BareRoot[4]{\frac{135}4}}\mathpunct{\cdot}\) Jetzt stelle ich die zweite Hypothese nicht auf gut Glück auf. Ich überlege vielmehr, dass sich eine Änderung von \(u\) bei mäßigen Zwischenzeiten immer viel stärker auf \(\log k\) als auf \(\tfrac12\log(r+r'')\) auswirkt, von \(\sec Q\) ganz zu schweigen. Ich rechne also den ganzen Fehler der ersten Hypothese \(\log k\) zu und gehe bei der zweiten Hypothese vom geänderten Wert von \(\log k\) aus, um daraus \(u\) mithilfe der Formeln
\[
\frac Ak=\sin\eta,\qquad u=k\cos\eta
\]
zu bestimmen. Nachdem ich auch diese Hypothese durchgerechnet und ihren Fehler ermittelt habe, bestimme ich den dritten Wert von \(u\) bereits durch Interpolation und habe damit schon gewonnenes Spiel. Sehen Sie hier die Anwendung auf das Beispiel aus meiner Abhandlung.

Die obige Formel ergibt
\[
\mathrm{I.}\quad u=0{,}15599.
\]

Aus der Zwischenzeit folgt \(\log\frac{t''-t}{m\BareRoot[4]{\frac{135}4}}=9{,}77932\). Aus diesem Wert \SourceLeaf{340}von \(u\), bei dem sich \(\log k=9{,}36029\) ergibt, folgt dafür hingegen \(9{,}65406\), also ein um 12526 Einheiten der fünften Stelle zu kleiner Wert. Für die zweite Hypothese nehme ich daher \(\log k=9{,}48555\), woraus folgt
\[
\mathrm{II.}\quad u=0{,}25561.
\]

Der Fehler dieser Hypothese beträgt \(+1598\), denn sie ergibt \(9{,}79530\). Daraus finde ich durch Interpolation als Verbesserung des zweiten \(u\) den Wert \(-0{,}01127\), also
\[
\mathrm{III.}\quad u=0{,}24434.
\]

Der Fehler dieser Hypothese beträgt \(+67\). Erneute Interpolation zwischen II und III ergibt also
\[
\mathrm{IV.}\quad u=0{,}24385,
\]
welcher Wert bis auf die letzte Einheit richtig ist. In meiner Abhandlung steht \(0{,}24388\), aber dort wird nach der exakten Formel gerechnet, die bei nur fünf Dezimalstellen weniger zuverlässig ist als die Näherungsformel. Sehen Sie hier die vollständige Durchführung des vierten Versuchs. Die roten\HistoricalNoteEight{Die im Original mit roter Tinte geschriebenen Zahlen sind hier unterstrichen.]} Zahlen sind bei allen Versuchen dieselben, und die Bedeutung der einzelnen Zahlen steht daneben.
\begingroup
\small
\setlength{\arraycolsep}{5pt}
\renewcommand{\arraystretch}{1.35}
\[
\begin{array}{rrr@{\quad}l}
0{,}24385&&&u\\
\underline{0{,}31365}&&&c\\
\underline{0{,}95443}&&&c''\\
\underline{9{,}74351}&\underline{9{,}91066}&\underline{9{,}22527}&bB,\ b''B'',\ A\\
9{,}74624&0{,}07856&9{,}38712&u+c,\ u+c'',\ u\\
\underline{0{,}17717}&\underline{0{,}05049}&9{,}91567&B\BareRoot{\frac{12}5},\ B''\BareRoot{\frac{12}5},\ \cos\eta\\
9{,}84812&9{,}74970&&\sin\theta,\ \sin\theta''\\
0{,}32905&&&s\\
0{,}30079&&&s''\\
0{,}28713&&&\text{\footnotemark}\\
0{,}61618&&&s+s''
\end{array}
\]
\[
\left.\begin{array}{r}9{,}47145\\0{,}30809\\-22\end{array}\right\}
9{,}77932\ \text{vollkommen genau}
\qquad
\left.\begin{array}{l}k\\\BareRoot{(s+s'')}\\-\BareRoot[5]{\sec Q}\end{array}\right\}\text{Ergebnis}
\]
\endgroup
\footnotetext{Aus der zweiten Spalte meiner Hilfstafel, wenn man in der ersten mit der Differenz der Logarithmen von \(s\) und \(s''\) nachschlägt.}

\pagebreak
\SourceLeaf{341}\noindent\hspace*{\parindent}.\ .\ .\ .\ .\ .\par
Ich bemerke noch, dass von den oben für die erste Näherung gemachten Voraussetzungen die erste bei mäßigen Zwischenzeiten völlig \textls[150]{legitim} ist. Insofern \(r\), \(r''\) ähnliche Funktionen von Größen sind, die sich nach dem Gesetz der Stetigkeit ändern, kann man das arithmetische Mittel \(\tfrac12(r+r'')\) näherungsweise als dieselbe Funktion der jeweiligen Mittelwerte ansehen. Die Ausdehnung auf \(k\) kann freilich in manchen Fällen allzu gewagt sein, doch ist der Versuch in jedem Fall unschädlich. Verzichtet man auf diese letzte Freiheit, so ist \(u\) durch die folgende Gleichung zu bestimmen
\[
\BareRoot[4]{(bb'')}\cdot\frac{t''-t}{m\BareRoot{18}}=\BareRoot{(uu+AA)}\,.\BareRoot[4]{\left\{\bigl(u+\tfrac12(c+c'')\bigr)^2+bb''BB''\right\}},
\]
die in entwickelter Form vom 6. Grad wäre, sich aber leicht durch Versuche lösen lässt. Dies wird wohl stets eine wirklich brauchbare Näherung ergeben. In unserem Beispiel folgt daraus \(u=0{,}2382\), wie Sie sehen, schon sehr nahe am wahren Wert. Am besten berechnet man wohl mithilfe dieser Gleichung nur eine oder zwei Ziffern von \(u\) und geht dann sogleich zu den genauen Versuchen über.
\par\medskip\noindent\makebox[\linewidth]{\rule{.18\linewidth}{.35pt}}\par
